\documentclass[11pt,a4paper]{article}
\usepackage[T1]{fontenc}
\usepackage[utf8]{inputenc}
\usepackage{amsmath,amssymb,amsthm}
\usepackage[margin=1in]{geometry}
\usepackage[colorlinks=true,linkcolor=blue,urlcolor=blue]{hyperref}
\theoremstyle{plain}
\newtheorem{theorem}{Theorem}
\newtheorem{lemma}[theorem]{Lemma}
\newtheorem{proposition}[theorem]{Proposition}
\newtheorem{corollary}[theorem]{Corollary}
\theoremstyle{definition}
\newtheorem{remark}[theorem]{Remark}

\title{The empirical recurrence for OEIS A297090, proved by transfer matrix}
\author{Adrian Perez Fontelles\\ \small Independent researcher}
\date{1 September 2026}
\begin{document}
\maketitle

\begin{abstract}
OEIS A297090 counts $n\times4$ arrays over $\{0,\dots,1\}$ subject to a condition imposed on
every CELL through its king-move neighbourhood, and it carries an empirical recurrence of
order $42$ contributed by R.~H.~Hardin. It is true, and it is decidable rather than
empirical. A cell's neighbourhood meets three consecutive rows, so the state of a transfer
matrix is a pair of consecutive rows and each step settles the middle one; an array is a
walk, and $a(n)$ is a walk count on $S=272$ vertices, which is $C$-finite. Whether the
conjectured recurrence annihilates it is then decided by a finite exact computation, with the
Cayley--Hamilton theorem bounding the work.
\end{abstract}

\noindent\small 2020 Mathematics Subject Classification. 05A15, 05B45, 15B36.\normalsize

\section{The conjecture}

OEIS A297090 is ``Number of nX4 0..1 arrays with no 1 adjacent to 3 king-move neighboring 1s.''. It has offset $1$ and begins
\[
16, 152, 1256, 11925, 113084, 1060921, 10038612, 95094637,\ \dots
\]
The entry states, as a conjecture contributed by its author and never marked settled:
\begin{quote}
Empirical: a(n) = 13*a(n-1) -32*a(n-2) +91*a(n-3) -1148*a(n-4) +1189*a(n-5) +432*a(n-6) +28862*a(n-7) +1885*a(n-8) -99288*a(n-9) -277764*a(n-10) -73602*a(n-11) +940696*a(n-12) +830603*a(n-13) +656028*a(n-14) -1842327*a(n-15) +670113*a(n-16) -4449019*a(n-17) +328722*a(n-18) -7184185*a(n-19) +5048924*a(n-20) -5904485*a(n-21) +17433108*a(n-22) +14043747*a(n-23) -6345609*a(n-24) -5376327*a(n-25) -14541923*a(n-26) +7951703*a(n-27) -46141383*a(n-28) +48123785*a(n-29) +10991042*a(n-30) +34162142*a(n-31) -12716751*a(n-32) +10255038*a(n-33) -6907418*a(n-34) -5580039*a(n-35) -2794607*a(n-36) +1015404*a(n-37) -7152*a(n-38) +25911*a(n-39) +630*a(n-40) -4704*a(n-41) -560*a(n-42)
\end{quote}
As of the ``Last modified'' line on the live entry (Dec 25 2017 09:48:52, revision 4) this is still
recorded as empirical, and nothing on the entry records it as proved.

\section{Cells, not blocks}

The king-move neighbours of a cell are the (at most) eight cells surrounding it. Every cell
carries the condition
\[
\text{every }1\text{ has }\#\{\text{neighbouring }1\text{s}\}\notin\{3\},
\]
and the count of neighbours enters it, so the boundary is not a detail that can be papered
over: a corner cell has three neighbours, an edge cell five, an interior cell eight, and a
condition such as ``a strict majority of its neighbours'' means different things at each. In
particular one cannot pad the array with a border of zeros --- that would give boundary cells
eight neighbours and change what is being counted. The outside is therefore carried
explicitly, as a sentinel row above the first and below the last and as a sentinel column
either side, and a neighbour that falls outside is simply absent from both counts.

Write an array as a sequence of rows $u_1,\dots,u_L$, each in
$\{0,\dots,1\}^{4}$. The neighbourhood of a cell in row $u_i$ meets $u_{i-1}$,
$u_i$ and $u_{i+1}$ and nothing else. Take as vertices the pairs
$(r,s)\in(\{0,\dots,1\}^{4}\cup\{\varnothing\})\times\{0,\dots,1\}^{4}$,
and put an edge $(r,s)\to(s,t)$ when row $s$ satisfies the condition at every one of its
$4$ cells, given $r$ above and $t$ below. An array with $L$ rows is then a walk of length
$L-1$ starting at $(\varnothing,u_1)$, each step settling one row, and the last row is
settled by the terminal weight, which evaluates it with $\varnothing$ below. Hence,
with $M$ the adjacency matrix, $\iota$ the starting vector and $\tau$ the terminal weight,
\[
a(n)\;=\;\,\iota^{\!\top}M^{\,n-1}\tau ,
\]
the exponent fixed by the entry's own shape and checked against its published terms. In
particular $a$ is $C$-finite.

\section{The proof}

\begin{lemma}\label{lem:crit}
Let $q(t)=t^{42}-\sum_i c_it^{42-i}$ be the characteristic polynomial of the
conjectured recurrence. Then the residual $a(n)-\sum_ic_ia(n-i)$ equals
$\iota^{\!\top}M^{\,j}q(M)\tau$ for the corresponding walk index $j$, so the
recurrence holds from some point on if and only if $\iota^{\!\top}M^{j}q(M)\tau=0$ for all
large $j$; and it is enough to examine $j<S$.
\end{lemma}

\begin{proof}
Factor $M^{j}$ out of $M^{j+42}-\sum_ic_iM^{j+42-i}$. For the bound, $M^{S}$
is an integer combination of $I,M,\dots,M^{S-1}$ by Cayley--Hamilton, so the numbers
$u_j=\iota^{\!\top}M^{j}q(M)\tau$ satisfy the monic linear recurrence given by the
characteristic polynomial of $M$; $S$ consecutive zeros force every later one, and the last
nonzero $u_j$ pins the threshold exactly.
\end{proof}

\begin{theorem}
\[
a(n)\;=\;13\,a(n-1) - 32\,a(n-2) + 91\,a(n-3) - 1148\,a(n-4) + 1189\,a(n-5) + 432\,a(n-6) + 28862\,a(n-7) + 1885\,a(n-8) - 99288\,a(n-9) - 277764\,a(n-10) - 73602\,a(n-11) + 940696\,a(n-12) + 830603\,a(n-13) + 656028\,a(n-14) - 1842327\,a(n-15) + 670113\,a(n-16) - 4449019\,a(n-17) + 328722\,a(n-18) - 7184185\,a(n-19) + 5048924\,a(n-20) - 5904485\,a(n-21) + 17433108\,a(n-22) + 14043747\,a(n-23) - 6345609\,a(n-24) - 5376327\,a(n-25) - 14541923\,a(n-26) + 7951703\,a(n-27) - 46141383\,a(n-28) + 48123785\,a(n-29) + 10991042\,a(n-30) + 34162142\,a(n-31) - 12716751\,a(n-32) + 10255038\,a(n-33) - 6907418\,a(n-34) - 5580039\,a(n-35) - 2794607\,a(n-36) + 1015404\,a(n-37) - 7152\,a(n-38) + 25911\,a(n-39) + 630\,a(n-40) - 4704\,a(n-41) - 560\,a(n-42)
\]
for every $n>42$.
\end{theorem}

\begin{proof}
The vector $w=q(M)\tau$ was formed by $42$ matrix--vector products in exact integer
arithmetic and $\iota^{\!\top}M^{j}w$ evaluated for $j=0,1,\dots$, again exactly. Those
integers vanish from the index corresponding to $n=42+1$ onwards, and the run was
continued until the working vector was identically zero, which settles all larger $j$ at once.
By Lemma~\ref{lem:crit} the recurrence holds for every $n>42$.
\end{proof}

\section{Verification}

Three checks, each independent of the algebra above.

First, the transfer model reproduces all $18$ terms the entry publishes, exactly, in
integer arithmetic. That check earned its place here. An early version of the construction
tested whether a neighbour lay in the current row by comparing the row objects rather
than their positions; when the row above happened to be identical to the current one, the
cell directly above was then skipped as though it were the cell itself. The counts agreed with
the entry for arrays of up to three rows and diverged at four, which is exactly the sort of
error that a check on a handful of small cases would have missed.

Second, the conjectured recurrence was evaluated directly on the entry's own published terms,
with no matrices involved, and holds wherever the proved range and the data overlap.

Third, the annihilation test of Lemma~\ref{lem:crit} was carried out in exact integer
arithmetic on the whole vector rather than on a sampled prefix.

\begin{thebibliography}{9}
\bibitem{oeis} The OEIS Foundation, \emph{The On-Line Encyclopedia of Integer
Sequences}, \url{https://oeis.org}, 2026. Sequence A297090.
\bibitem{stanley} R.~P.~Stanley, \emph{Enumerative Combinatorics}, Volume 1, 2nd ed.,
Cambridge University Press, 2011. (Transfer-matrix method, Section 4.7.)
\bibitem{fs} P.~Flajolet and R.~Sedgewick, \emph{Analytic Combinatorics}, Cambridge
University Press, 2009.
\bibitem{horn} R.~A.~Horn and C.~R.~Johnson, \emph{Matrix Analysis}, 2nd ed., Cambridge
University Press, 2013.
\end{thebibliography}

\end{document}
